\section{Chapitre 3 : Architecture et Conception Systémique}
\addcontentsline{toc}{section}{Chapitre 3 : Architecture et Conception Systémique}

La conception architecturale de la plateforme Fraudly repose sur l'adoption délibérée d'un modèle en microservices, une décision stratégique motivée par la nécessité impérieuse de découpler les domaines fonctionnels hétérogènes qui composent ce système éducatif complexe. Contrairement à une approche monolithique conventionnelle où l'intégralité de la logique métier réside au sein d'un processus unifié, l'architecture en microservices prône la ségrégation de l'application en une myriade de services logiciels indépendants, autonomes et circonscrits à un périmètre métier précis. Cette décentralisation structurelle offre des avantages considérables en matière de résilience globale et d'évolutivité dynamique. En effet, lors des périodes critiques de forte affluence, telles que les sessions d'examens simultanées de fin de semestre, les composants spécifiquement sollicités, à l'instar du service de télésurveillance ou du service de gestion des évaluations, peuvent être répliqués horizontalement sur des grappes de serveurs pour absorber la charge de calcul, sans pour autant nécessiter la redondance inutile des services périphériques moins sollicités. De surcroît, cette architecture polyglotte permet d'attribuer à chaque microservice l'écosystème technologique le plus idoine à la résolution de sa problématique intrinsèque. Ainsi, les processus de gestion transactionnelle stricte et d'orchestration des données relationnelles ont été confiés à l'écosystème Java par le biais du cadriciel Spring Boot, reconnu pour sa robustesse opérationnelle et sa rigueur structurelle, tandis que les opérations hautement calculatoires inhérentes à l'intelligence artificielle générative et au traitement asynchrone des flux vidéo ont été déléguées à un environnement d'exécution Python, tirant parti de la puissance asynchrone du cadriciel FastAPI et de l'incomparable richesse de son écosystème scientifique dédié à l'apprentissage profond.

Le socle transactionnel de l'architecture est structuré autour de plusieurs services indépendants développés sous Spring Boot, communiquant via des interfaces de programmation applicative de type Representational State Transfer pour les requêtes synchrones. Le service d'authentification centralise la gestion des identités, s'assurant de la vérification des accréditations des utilisateurs, de la gestion du cycle de vie des mots de passe et de la génération de jetons d'accès chiffrés, garantissant ainsi l'étanchéité sécuritaire aux frontières du système. Le service de gestion des cours orchestre la structuration arborescente du contenu pédagogique, administrant la hiérarchie complexe entre les programmes de formation, les modules constitutifs et les ressources granulaires multimédias qui les composent, tout en régissant les droits d'accès inhérents aux inscriptions des étudiants. Le service d'évaluation endosse la responsabilité critique de l'orchestration des examens, coordonnant la distribution des questions chronométrées, capturant de manière persistante les tentatives de réponses des apprenants et appliquant les règles déterministes de notation pour les questionnaires à choix multiples. En parallèle, le service d'analyse de l'apprentissage agrège continuellement les signaux d'engagement et les résultats quantitatifs des étudiants pour forger des profils d'apprentissage détaillés, stockant méticuleusement les métriques de réussite, les compétences non maîtrisées et les statistiques de progression au sein d'une base de données relationnelle dédiée. Enfin, le service de télésurveillance agit comme l'organe de réception des événements comportementaux en temps réel, consignant les anomalies d'interaction signalées par l'interface utilisateur et centralisant la génération des rapports d'intégrité académique à l'issue de chaque session de contrôle continu.

Afin de pallier les limitations inhérentes à la communication synchrone entre ces multiples composants fortement découplés, l'architecture intègre un bus de messages distribué, matérialisé par la technologie Apache Kafka. Ce courtier de messages s'avère particulièrement fondamental pour orchestrer les processus asynchrones dont la latence d'exécution est imprévisible ou significativement élevée, tels que la sollicitation des grands modèles de langage pour la notation sémantique. Lorsqu'un étudiant soumet une réponse textuelle élaborée à une question ouverte, le service d'évaluation Java ne bloque pas l'interface en attendant la réponse de l'intelligence artificielle, mais se contente de publier un événement de demande de correction sur un canal de distribution Kafka spécifique. Ce message est ultérieurement consommé, à son propre rythme, par le service d'intelligence artificielle hébergé sous Python, qui interrogera le modèle de langage en lui soumettant le triptyque constitué de la question, du barème de notation et de la proposition de l'étudiant. Une fois le raisonnement évaluatif achevé et la note sémantique attribuée, l'agent artificiel publie à son tour un événement de résultat sur un second canal Kafka. Le service d'évaluation d'origine, agissant cette fois-ci en tant que consommateur, réceptionne ce résultat asynchrone, met à jour le registre permanent des notes dans la base de données relationnelle PostgreSQL et consolide le score final de la tentative d'examen. Ce paradigme de communication orientée événements par abonnement et publication assure une résilience remarquable du système : en cas d'indisponibilité temporaire du fournisseur d'intelligence artificielle ou de pic d'affluence extrême, les demandes de correction s'accumulent de manière sécurisée dans la file d'attente Kafka, éliminant ainsi le risque dramatique de perte de données d'examen et préservant l'expérience fluide de l'utilisateur final qui peut poursuivre sa navigation sans interruption.

La stratégie de persistance des données au sein de Fraudly obéit au principe de la polyglottie des bases de données, stipulant qu'aucune technologie de stockage unique ne saurait répondre optimalement aux exigences contradictoires de tous les domaines métiers. Pour les données hautement structurées exigeant le respect inaltérable des propriétés d'atomicité, de cohérence, d'isolation et de durabilité, telles que les transactions d'authentification, les configurations des cours, les barèmes de notation et l'historique financier des inscriptions, le système s'appuie sur le moteur de base de données relationnelle PostgreSQL. Chaque microservice métier dispose de son propre schéma de base de données physiquement ou logiquement isolé, empêchant ainsi tout couplage par la donnée et forçant l'accès exclusif aux informations via les interfaces de programmation applicatives autorisées. En revanche, le module d'intelligence artificielle en charge de la fonctionnalité de tuteur virtuel requiert un mécanisme de persistance fondamentalement différent, optimisé pour la recherche de similarité sémantique plutôt que pour l'appariement syntaxique exact. À cet effet, l'architecture déploie une base de données vectorielle spécialisée, ChromaDB. Lorsqu'un enseignant téléverse un nouveau document de cours en format texte ou extrait via la reconnaissance optique de caractères, le contenu est segmenté en paragraphes cohérents puis transformé par un modèle d'enchâssement en une myriade de vecteurs numériques à haute dimensionnalité, qui sont alors indexés au sein de ChromaDB. Cette topologie structurelle mixte confère au système Fraudly la fondation nécessaire pour soutenir à la fois la rigueur exigée par un système d'information académique certifiant et l'agilité indispensable au déploiement d'agents cognitifs de pointe.
